\chapter{Gráfépítés, klaszterezés és perzisztencia}

\section{A jelenlegi gráfépítés szereposztása}

A rendszer alapgráfja meccsekből épül, de a jelenlegi architektúrában már nem a Python/PyVis export a fő gráfépítő út. A korai Python pipeline fontos prototípus- és kompatibilitási réteg maradt, de a dolgozat szempontjából az elsődleges, kutatható gráfállapotot a Rust réteg állítja elő. A Python oldal megőrzi az adatgyűjtés, gyors prototipizálás, régi co-presence export és kompatibilitási gráfgenerálás szerepét; az SQLite adatbázis tartja a normalizált játékosprofilokat, score-okat és a perzisztált klasztertagságokat; a Rust motor építi fel a futásidejű \textit{GraphState}-et a meccs-JSON állományokból és az adatbázisból; a Graph V2 Rust pipeline állítja elő a datasethez kötött, cache-elhető gráf-artifaktumokat, amelyeket a frontend előkészített bináris és JSON állományokként fogyaszt.

Ez a pontosítás fontos módszertani korrekció: a fejezet nem állíthatja azt, hogy a Python pipeline a klaszterezés és gráfépítés fő útja. A Python út történeti és debug-szerepe megmarad, de a \textit{flexset} és \textit{soloq} összehasonlítás, az asszortativitás, a központisági elemzés és a Graph V2 vizualizáció már Rust-központú gráfmodellre támaszkodik.

\section{Kapcsolataggregáció a Rust gráfmodellben}

A Rust gráfmodell minden játékospárra külön őrzi a szövetségesi és ellenfél-oldali bizonyítékot. Ha a mérkőzések halmazát $M$, egy mérkőzés két csapatát pedig $T_{m,1}$ és $T_{m,2}$ jelöli, akkor egy játékospárra a két fő súly a következőképpen írható le:

\begin{equation}
w_{ally}(u,v)=\sum_{m \in M}\mathbf{1}[\exists i \in \{1,2\}: u \in T_{m,i} \land v \in T_{m,i}]
\end{equation}

\begin{equation}
w_{enemy}(u,v)=\sum_{m \in M}\mathbf{1}[\exists i \neq j: u \in T_{m,i} \land v \in T_{m,j}]
\end{equation}

Ebből a rendszer három alapmezőt képez: az \textit{allyWeight} azt mutatja, hányszor szerepelt a két játékos ugyanazon az oldalon; az \textit{enemyWeight} azt, hányszor kerültek egymással szembe; a \textit{totalMatches} a kettő összege. A domináns reláció \textit{ally}, ha \textit{allyWeight} $\geq$ \textit{enemyWeight}, különben \textit{enemy}. A dolgozat fő olvasata szempontjából ez nem azt jelenti, hogy az enemy él társadalmi ellenségességet bizonyít: a Rust modell ismétlődő meccsbeli együtt- és szembekerülési bizonyítékot tárol, ezért a Graph V2 eredményei elsősorban asszociatív játékosgráfként, nem szó szerinti barátsági vagy ellenségességi hálóként értelmezendők.

\section{Szűrés, Graph V2 és bounded ally csoportok}

A zajcsökkentés alapvető tervezési döntés. A Graph V2 láthatósági küszöbe \textit{totalMatches} $\geq 2$: egyetlen közös mérkőzés még nem elég ahhoz, hogy egy kapcsolat a populációs renderben vagy a kutatási artifact-rétegben stabil élként jelenjen meg. Magasabb küszöbnél a háló tisztább, de sok gyengébb szerkezet eltűnik; alacsonyabb küszöbnél nő a véletlenszerű queue-találkozások zaja. Ez a küszöb nem abszolút igazság, hanem interpretálhatósági kompromisszum.

A Graph V2 klaszterezés fő iránya a Rust oldali, determinisztikus \textit{bounded-ally-groups-v2} logika, amelynek alapja az \textit{allyWeight} $\geq 2$ típusú ismétlődő szövetségesi kapcsolat. Az enemy-only élek nem olvasztanak össze klasztereket, mert ezzel a közösségértelmezés túl agresszívan és gyengén védhetően terjedne szét. Az enemy információ megmarad az élmodellben és megjeleníthető külön rétegként, de nem ez határozza meg az ally-csoportok alapját. A vizuális klaszterek mérete szándékosan korlátozott, legfeljebb körülbelül húsz játékosra, megakadályozva, hogy hosszú tranzitív teammate-láncok egyetlen, hamisan értelmezhető óriásklaszterré olvadjanak.

\section{Graph V2 artifaktumok mint igazságforrás}

A jelenlegi gráfépítés egyik legfontosabb különbsége a korai Python/PyVis exporthoz képest az, hogy a böngésző már nem kap egy nagy, önmagában layoutolt HTML-gráfot. A Rust Graph V2 pipeline előkészített, datasethez kötött cache-artifaktumokat ír ki: a \textit{manifest.json} rögzíti a datasetet, a gráfépítő verziót, a klaszterezési és layout verziót, a csúcs- és élszámokat; a \textit{node\_positions.f32}, \textit{edge\_pairs.u32} és társaik a Three.js/WebGL réteg bináris bemeneti állományai; a \textit{cluster\_meta.json} a klaszterazonosítókat és tagságokat tartalmazza; a \textit{summary.md} reprodukálhatósági és elemzési kontextust rögzít: kimondja, milyen adathalmazon, milyen küszöbökkel és milyen konstrukciós szabállyal készült az adott gráf. A frontend ezeket az artifaktumokat fogyasztja, nem önálló gráfépítő igazságforrásként működik. A Graph V2 artifact-réteg tekintendő a kutatási vizualizáció fő igazságforrásának; a Python HTML-export legfeljebb történeti, debug- vagy kompatibilitási nézet.

\section{Klaszterek és a Python pipeline státusza}

A jelenlegi rendszerben a klaszterek önálló, lekérdezhető adatbázis-objektumok; a frontendnek nem kell minden alkalommal újragenerálnia a klaszterlogikát, és a Rust futtatókörnyezet és a Python pipeline ugyanazt a perzisztens modellt tudja használni. A sémáról és a két klasztercsaládról az 5. fejezet szól részletesen; itt elegendő rögzíteni, hogy a gráfépítési szerepkiosztás szempontjából a \textit{python\_population} klaszterek történeti, populációs és debug-szerepet töltenek be, míg a \textit{rust\_pathfinding} klaszterek a jelenlegi dolgozati főútvonalhoz igazodnak: a datasetváltáshoz, az útvonalkereséshez, a Graph V2 artifactokhoz és az analitikai modulokhoz.

\diagramfigure{A Pathfinder osztálydiagramja - a gráf- és klaszterobjektumok tárolási modellje}{path finder class diagram.pdf}

A \textit{backend/build\_graph.py} és a hasonló Python gráfépítő modulok megmutatják, honnan indult a projekt: gyors NetworkX/PyVis alapú prototípusként, amely súlyozott co-presence gráfot épített, HTML/JSON exportokat készített, majd fokozatosan adatbázis-perzisztenciával egészült ki \cite{networkx_greedy,pyvis}. A jelenlegi fejezetben ezt már nem fő kutatási pipeline-ként, hanem visszaminősített, legacy-kompatibilitási útként kezeljük; a fő empirikus értelmezés a \textit{flexset} és \textit{soloq} Rust-gráfjaira támaszkodik.

\section{Ország- és régióelemzés mint korábbi kísérleti ág}

A projekt egy korábbi szakaszában a klaszterekből nemcsak gráfstruktúrát, hanem földrajzi következtetéseket is szerettünk volna kinyerni. A \textit{fetch\_clusters.py}, \textit{assign\_countries.py} és \textit{conclude.py} modulok klaszternevekből kísérleti országcímkézést futtattak és aggregált ország-alapú összegzést készítettek. Az elképzelés vonzó volt, mert a puszta gráfstruktúrán túl egyfajta társas magyarázó réteget ígért; jól mutatja azt is, hogy a projekt nem egyetlen előre rögzített hipotézis mentén fejlődött, hanem több analitikai irányt próbált ki.

A modul végül nem vált a dolgozat fő empirikus pillérévé. Az ország-hozzárendelés soha nem közvetlen Riot-mezőből származott, hanem névmintázatokból következtetett, módszertanilag bizonytalan címke volt; az ellenőrzött \textit{playersrefined.db} pillanatképben a \textit{country} mező érdemi kitöltöttség nélkül maradt. Az ág önálló fejezetet nem indokol, de a kódnyomok megőrzésre kerültek, és a rendszer gondolkodásmódját jól dokumentálják.

\diagramfigure{A korai ország-klaszter pipeline szekvenciadiagramja - a frontend, a backend és a Python worker közötti folyamat az ország-hozzárendelési kísérletben}{old graph sequence diagram.pdf}
